% SilentEHR — ACL-style manuscript (compile with pdflatex + acl.sty)
\documentclass[11pt]{article}

\usepackage[review]{acl}
\usepackage{times}
\usepackage{latexsym}
\usepackage{amsmath,amssymb}
\usepackage{graphicx}
\usepackage{booktabs}
\usepackage{placeins}
\usepackage{multirow}
\usepackage{array}
\usepackage{xcolor}

\title{SilentEHR: Auditing Closed-World Extrapolation\\in Frozen Patient Representations}

\author{
  Anonymous ACL submission
}

\begin{document}
\maketitle


\begin{abstract}
Binary clinical verifiers are commonly trained on records that explicitly
support or refute a proposition, then applied to records that do not mention it.
We study closed-world extrapolation (CWE): mapping non-observation to Refuted
rather than Unknown. SilentEHR provides 11{,}466 controlled
patient--proposition pairs with held-out background, template, and concept
protocols. A complementary MIMIC-IV analysis separates laboratory state from
textual observation. Across seven frozen encoders, Unmentioned-to-Refuted rates
reach 0.99, although six models show no geometric collapse toward Refuted. On
two human-consensus MIMIC concepts, dense encoders predict absence for 41--55\%
of structured-positive cases labeled unmentioned by Qwen3-14B; observation
detection remains near chance. A 300-case dual-annotator audit finds high human
agreement ($\kappa=.80$) but imperfect LLM labels (cohort-weighted $O^-$ precision .60,
95\% CI [.42,.78]). Artifact controls preserve the directional result.
Explicit Unknown supervision reduces CWE, whereas confidence- and
Mahalanobis-based abstention do not.
\end{abstract}

\section{Introduction}
Clinical text often supports patient-level verification questions: whether a
condition is present, absent, possible, historical, or unresolved. Natural
language inference (NLI) formalizes related distinctions as entailment,
contradiction, and neutral \citep{bowman-etal-2015-large,williams-etal-2018-broad},
and clinical variants evaluate inference over patient histories and trial
reports \citep{romanov-shivade-2018-lessons,jullien-etal-2023-nli4ct}.
Clinical assertion systems likewise distinguish explicit presence, possibility,
and absence for mentioned entities \citep{van-aken-etal-2021-assertion}.
These settings assume that a model is trained to recognize the relevant
third state. A common lightweight deployment pattern pairs frozen representations
with a binary downstream head trained only on explicit positive and negative
evidence. At deployment, the record may simply not mention the proposition.

Silence is not equivalent to negation. Clinical notes are selective summaries,
and prior work has treated explicit evidence of absence as a distinct extraction
problem \citep{carter-etal-2016-evidence}. When a proposition is unmentioned, a
safe verifier should return Unknown or abstain rather than infer Refuted. Yet a
binary classifier has no supervised Unknown region and must extrapolate. Average
accuracy on observed examples does not reveal the direction of this
extrapolation, whether it is expressed with high confidence, or whether it is
already encoded in the geometry of the frozen representation.

We study this failure mode, which we call \emph{closed-world extrapolation}
(CWE). SilentEHR creates matched counterfactual profiles that differ only in
epistemic state and evaluates held-out backgrounds, assertion templates, and
clinical concepts. A complementary MIMIC-IV analysis separates structured
laboratory state from textual observation in deidentified discharge summaries
\citep{johnson-etal-2023-mimiciv,johnson-etal-2023-mimicnote}. Together, the two
layers provide internal control and real-note relevance: SilentEHR isolates
where a binary decision boundary sends non-observation, while MIMIC-IV tests a
related documentation-mismatch consequence under imperfect observation labels.

The results reveal a reproducible separation between direction, confidence, and
geometry. Several dense encoders send more than 90\% of Unmentioned profiles to
Refuted, often confidently, although their Unmentioned embeddings do not
collapse toward the Refuted centroid. The behavior persists after controlling
text length, evidence position, and direct lexical assertions, but varies across
encoders, concepts, and fine-grained Unknown states. On human-consensus MIMIC-IV
concepts, dense encoders also predict clinical absence for 41--55\% of
structured-positive cases labeled textually unmentioned, while confidence does
not reliably distinguish observed from unmentioned notes.

Our contributions are threefold. First, we introduce SilentEHR, an 11{,}466-pair
controlled audit benchmark with matched epistemic states and multiple transfer
protocols. Second, we propose diagnostics that separately measure extrapolation
direction, confidence extremity, and representation geometry, and show that
these quantities need not agree. Third, we combine structured-grounded MIMIC-IV
validation, a 300-case dual-annotator audit, and mitigation experiments showing
that explicit Unknown supervision is substantially more effective than
confidence- or distance-based abstention.

\section{Datasets}
Table~\ref{tab:datasets} summarizes the two evaluation layers. SilentEHR
isolates CWE under controlled censoring; MIMIC-IV tests a related
documentation-mismatch consequence in real notes.


\begin{table}[t]
  \centering
  \footnotesize
  \renewcommand{\arraystretch}{1.06}
  \caption{Evaluation layers.}
  \label{tab:datasets}
  \begin{tabular*}{\columnwidth}{@{\extracolsep{\fill}}lcc@{}}
    \toprule
    & \textbf{SilentEHR} & \textbf{MIMIC-IV} \\
    \midrule
    Contexts & 80 synthetic & 13.8k admissions \\
    Concepts & 16 & 3 laboratory \\
    Pairs & 11,466 & 38,043 \\
    States & S/R/U/NA/Unc & $C^{\pm}O^{\pm}$ \\
    Observation & controlled & Qwen3-14B \\
    Split & background--concept & subject \\
    \bottomrule
  \end{tabular*}
\end{table}


\subsection{SilentEHR}
Sixteen concepts are crossed with 80 target-neutral backgrounds. After removing
six contaminated background--concept pairs, each of 1{,}274 pairs yields two
Supported, two Refuted, two Not-assessed, two Uncertain, and one Unmentioned
variant, totaling 11{,}466 examples. IID assignment is performed at the
background--concept-family level using a stable hash; all state and template
realizations of one family remain in the same fold. The primary split trains on
56 backgrounds, validates on 12, and tests on 12; template-held-out and
concept-held-out protocols provide complementary tests. Concepts span symptoms,
allergies, medications, disease states, and safety conditions.

Three regenerated tiers control common artifacts. Length matching adds a
concept-neutral sentence to Unmentioned profiles, reducing the clean-tier
length gap from 5.8 words to less than one word. Random-position evidence
removes the fixed final-sentence cue. Indirect evidence replaces explicit
assertions with laboratory values or medication events. The full corpus and
binary head are regenerated and retrained for each tier.

\paragraph{Matched-family example.}
Each background--concept pair shares one synthetic background; only the final
state sentence changes (Appendix~\ref{app:matched_family}).
Table~\ref{tab:matched_family_compact} shows the \texttt{penicillin\_allergy}
family on background~0 (\texttt{p01440--p01448} in \texttt{controlled\_full.jsonl}).

\begin{table}[t]
  \centering
  \footnotesize
  \renewcommand{\arraystretch}{1.08}
  \caption{Compact matched-family example (background~0, penicillin allergy).}
  \label{tab:matched_family_compact}
  \begin{tabular*}{\columnwidth}{@{\extracolsep{\fill}}lp{0.62\columnwidth}@{}}
    \toprule
    \textbf{State} & \textbf{Example realization (verbatim)} \\
    \midrule
    Proposition & ``The patient has a penicillin allergy.'' \\
    Supported & ``A penicillin allergy is documented.'' \\
    Refuted & ``The record states no known penicillin allergy.'' \\
    Not assessed & ``Penicillin allergy status was not documented.'' \\
    Uncertain & ``The reported penicillin allergy could not be confirmed.'' \\
    Unmentioned & Shared background only; no penicillin sentence appended. \\
    \bottomrule
  \end{tabular*}
\end{table}

\subsection{MIMIC-IV}
We use MIMIC-IV and MIMIC-IV-Note
\citep{johnson-etal-2023-mimiciv,johnson-etal-2023-mimicnote}, pairing discharge
summaries with the final laboratory measurement within 24 hours of discharge: hemoglobin $<10$ g/dL, platelets
$<100\times10^9$/L, or potassium $>5.5$ mmol/L. Qwen3-14B assigns one of seven observation statuses from the full note:
current positive, current negative, unmentioned, admission-only, historical,
uncertain, or temporally conflicting (Appendix~\ref{app:observation_labeling}).

Two experiments use these labels differently.
For the \emph{direct textual S/R/U bridge} (Appendix~\ref{app:real_sru}),
\emph{unmentioned}, \emph{admission-only}, \emph{historical-only},
\emph{uncertain}, and \emph{temporal-conflict} all map to Unknown.
For the separate \emph{structured-state} $C\times O$ analysis
(Appendix~\ref{app:mimic_controls}), uncertain, historical-only,
admission-only, and temporal-conflict cases are \emph{excluded} from binary
classifier training and evaluation to avoid conflating observation ambiguity
with laboratory-state prediction; only clean observed-positive/negative and
unmentioned cells enter that protocol.
Full notes are encoded by mean pooling overlapping 4,800-character chunks with
400-character overlap.

Two graduate-student annotators independently label 300 stratified anemia and
thrombocytopenia cases under a blind protocol (Appendix~\ref{app:annotation_protocol}):
150 anemia + 150 thrombocytopenia, using the same seven-class schema and quoting
evidence for every non-\emph{unmentioned} label. We report human--human
agreement, human--LLM agreement, dual-annotator consensus labels
($n=305$ after a 50-pair non-laboratory expansion; 255 from the original audit),
binary mentioned/unmentioned metrics, and the precision of LLM-predicted $O^-$.
Because the audit is stratified, we additionally reweight precision to the MIMIC
test distribution.

\section{Methods}
\subsection{Closed-world extrapolation}
For record $x$ and proposition $h$, the target is
$y(x,h)\in\{S,R,U\}$. Training uses only Supported and Refuted cases.
UNR is the fraction of Unmentioned cases predicted Refuted; HCUNR@0.9 applies a
0.9 confidence threshold. Detection AUROC uses low maximum class probability as
a non-observation score. Matched ECI, $1-d(U,R)/d(S,R)$ in PCA--Ledoit--Wolf Mahalanobis space,
measures whether Unmentioned embeddings collapse toward Refuted. Negative ECI
means that high UNR cannot be explained by global proximity to Refuted. We also
project examples onto the learned logistic-regression axis to distinguish
global geometry from supervised extrapolation along the decision direction.

\subsection{Encoders and classifier}
We evaluate TF--IDF; Bio\_ClinicalBERT
\citep{alsentzer-etal-2019-clinicalbert}; BioLORD
\citep{remy-etal-2022-biolord}; MedCPT \citep{jin-etal-2023-medcpt};
Qwen3-Embedding-0.6B \citep{zhang-etal-2025-qwen3embedding}; Jina v5
\citep{akram-etal-2026-jinav5}; and med-gte-hybrid
\citep{reimers-gurevych-2019-sentence,medaihs-2025-medgte}, all with frozen
weights. Pair features are
$[e_x;e_h;|e_x-e_h|;e_x\odot e_h]$. The primary head is class-balanced $\ell_2$-regularized logistic regression
on standardized features, fitted only to Supported and Refuted training pairs.
Robustness tests vary regularization strength, PCA dimension, $k$NN and
nearest-centroid heads, and temperature scaling. Neural encoder weights remain
frozen in every experiment.

\subsection{Evaluation protocol}
All preprocessing, standardization, dimensionality reduction, and decision
thresholds are fitted on the training or validation partition only. The primary
SilentEHR result uses one fixed background split for direct model comparison.
We then redraw ten background splits and four disjoint concept folds to test
partition sensitivity, retraining the classifier on cached embeddings.
Background--concept bootstrap intervals (2{,}000 replicates, resampling unit =
background$\times$concept family) preserve the matched counterfactual
unit; matched variants are never treated as independent samples.
Known-state accuracy is always reported with unknown-state behavior so
that high UNR cannot be attributed to failure on explicit evidence.
We report a 0.5 UNR reference as a \emph{directional} threshold (whether
silence is mapped toward Refuted), not as a null model of random
three-way classification performance.

Confidence-based detection uses maximum class probability, a common baseline
for error and out-of-distribution detection \citep{hendrycks-gimpel-2017-baseline},
without access to Unknown examples. Confidence abstention selects a validation threshold
that retains 95\% of known cases. Mahalanobis detection follows class-conditional distance-based OOD scoring
\citep{lee-etal-2018-mahalanobis}, using PCA and a Ledoit--Wolf covariance
estimate fitted to known training representations. These
settings deliberately test whether a conventional binary pipeline can recognize
non-observation after training, rather than giving the detector labeled Unknown
examples.

\subsection{MIMIC state--observation evaluation}
Clinical state $C$ comes from laboratories and observation $O$ from Qwen3-14B.
This structured $C\times O$ protocol (Appendix~\ref{app:mimic_controls}) is
separate from the direct textual S/R/U bridge in Appendix~\ref{app:real_sru}.
Training uses $C^+O^+$ and $C^-O^+$ on clean observed-positive/negative cells;
evaluation adds LLM-labeled $C^+O^-$ and $C^-O^-$ where the concept is
\emph{unmentioned} near discharge. Uncertain, historical-only, admission-only,
and temporal-conflict labels are excluded here (but map to Unknown in the S/R/U
bridge). Subject-level 70/30 splits prevent leakage, and neural models cap each
concept--state cell at 500 cases. The main result pools the two audited concepts
($n=75$ positive $O^-$ cases); hyperkalemia is exploratory.

UPFNR is $1-\mathrm{Recall}(C^+\mid O^-)$, HCWE is the high-confidence wrong
negative rate, and the documentation gap compares recall on observed and
LLM-labeled unmentioned cases. Observation-detection AUROC uses model
confidence to separate $O^-$ from $O^+$. These quantities measure
documentation-sensitive behavior rather than hidden-laboratory recovery.

A proposition-only classifier and a concept-prior baseline test whether low
UPFNR can be obtained by ignoring patient text. Label-noise sensitivity uses
the cohort-weighted $O^-$ precision $q=.60$ as a conservative setting. We
report both a missing-at-random deletion scenario and worst-case
partial-identification bounds that allow incorrect $O^-$ labels to concentrate
on model errors or correct predictions.

MIMIC confidence intervals use subject-level bootstrap sampling (2{,}000
replicates) and pool concept results by the number of structured-positive $O^-$
cases. Real-note S/R/U gold uses Wilson 95\% intervals for binomial rates on
pooled out-of-fold predictions. The human-consensus headline excludes
hyperkalemia because its positive $O^-$ test cell contains only five cases;
the three-concept result remains in the exploratory table.
For the 300 audited rows, we additionally recompute UPFNR using each annotator,
their consensus, union, and intersection definitions of $O^-$. These small
subsets are treated as directional sensitivity analyses rather than precise
effect-size estimates.

\subsection{Mitigation}
We compare temperature scaling \citep{guo-etal-2017-calibration}, confidence
and Mahalanobis abstention,
and three-way training with 10, 50, or 100 Unknown semantic groups over ten
seeds. An oracle observation gate provides an upper bound by abstaining on
Unknown while preserving binary predictions for known states.

\section{Results}
\subsection{Controlled benchmark}
Table~\ref{tab:controlled} shows strong but non-universal negative
extrapolation. Qwen3-Emb and BioLORD reach Unmentioned UNR .94 [.90,.98] and
.99 [.98,1.00] (family-bootstrap 95\% CI on the new-background split) while
maintaining known accuracy above .97. Not-assessed behaves similarly, whereas
Uncertain is encoder-dependent (Jina: .18). Med-gte is a clean-tier
counterexample (.38), although its rate rises under position and indirect
controls. Detection AUROC is at or below chance for several high-UNR encoders:
Jina reaches .498, BioLORD .184, and MedCPT .300. Only BioLORD has positive
ECI; Qwen and Jina remain negative despite strong directional failure.
Decision-axis logits further show that Jina maps Uncertain toward Supported
while leaving Unmentioned strongly negative. Thus Unknown states are not
monolithic, and prediction direction, confidence, and global geometry do not
coincide.


\begin{table*}[t]
  \centering
  \footnotesize
  \renewcommand{\arraystretch}{1.07}
  \caption{SilentEHR new-background results. U, NA, and Unc denote
  Unmentioned, Not-assessed, and Uncertain. Direction columns report the
  fraction predicted Refuted; HC uses a 0.9 confidence threshold. Det. is
  Unmentioned detection AUROC and ECI is the matched collapse index.}
  \label{tab:controlled}
  \begin{tabular*}{\textwidth}{@{\extracolsep{\fill}}lrrrrrrrrr@{}}
    \toprule
    & & \multicolumn{3}{c}{\textbf{Predicted Refuted}} &
    \multicolumn{3}{c}{\textbf{HC Refuted}} & & \\
    \cmidrule(lr){3-5}\cmidrule(lr){6-8}
    \textbf{Encoder} & \textbf{Known} &
    \textbf{U} & \textbf{NA} & \textbf{Unc} &
    \textbf{U} & \textbf{NA} & \textbf{Unc} &
    \textbf{Det.} & \textbf{ECI} \\
    \midrule
    TF--IDF      & 1.00 & .69 & .69 & .69 & .45 & .23 & .46 & .999 & $-.29$ \\
    ClinicalBERT & .98 & .60 & .32 & .55 & .27 & .20 & .24 & .866 & $-1.16$ \\
    Qwen3-Emb    & .98 & .94 & .91 & .56 & .88 & .86 & .48 & .657 & $-.15$ \\
    Jina v5      & .96 & .91 & .72 & .18 & .81 & .46 & .07 & .498 & $-.94$ \\
    BioLORD      & .99 & .99 & .77 & .65 & .99 & .64 & .56 & .184 & $.10$ \\
    MedCPT       & .89 & .79 & .19 & .43 & .72 & .11 & .31 & .300 & $-3.95$ \\
    med-gte      & 1.00 & .38 & .47 & .53 & .16 & .26 & .25 & .747 & $-2.16$ \\
    \bottomrule
  \end{tabular*}
\end{table*}


\subsection{Stability and heterogeneity}
Across ten background splits, Qwen, Jina, BioLORD, and TF--IDF exceed UNR .5
in every split, with means of .90, .85, .97, and .71; med-gte remains below .5
in all splits (.34). Four concept folds show the same family-level pattern.
Qwen, Jina, and BioLORD also predict Refuted for Unmentioned examples in all 12
training concepts, whereas med-gte does so in four. Its apparent safety is
concentrated in allergy and medication concepts, while disease concepts remain
high (.92), and disappears under indirect evidence. Full split and concept
results are in the appendix.

The concept-fold results also remain high for the strongest encoders, indicating
that the binary direction transfers beyond the concepts used to fit the head.
However, the med-gte counterexample cautions against treating CWE as universal:
encoder objectives and lexical realization materially change whether silence is
mapped to Refuted, Supported, or a low-margin region.

\subsection{Artifact controls}
Table~\ref{tab:artifact} summarizes the three artifact-resistant tiers.
Length matching leaves Qwen, Jina, and BioLORD UNR at .88, .80, and .99.
Randomized position also preserves the strongest failures. Indirect evidence reduces Qwen and BioLORD modestly but leaves Jina at .91;
TF--IDF rises to 1.00, consistent with strong lexical reliance. Known-state
accuracy remains high for the headline dense encoders in the length-matched and
indirect tiers. CWE therefore cannot be explained by sentence count, fixed
position, or direct target wording alone.


\begin{table}[t]
  \centering
  \scriptsize
  \setlength{\tabcolsep}{2.5pt}
  \renewcommand{\arraystretch}{1.05}
  \caption{Unmentioned-to-Refuted rate under artifact controls.}
  \label{tab:artifact}
  \begin{tabular*}{\columnwidth}{@{\extracolsep{\fill}}lrrrr@{}}
    \toprule
    \textbf{Encoder} & \textbf{Clean} & \textbf{Length} &
    \textbf{Position} & \textbf{Indirect} \\
    \midrule
    TF--IDF      & .688 & .688 & .688 & 1.000 \\
    ClinicalBERT & .597 & .764 & .153 & .833 \\
    Qwen3-Emb    & .944 & .882 & .944 & .785 \\
    Jina v5      & .910 & .799 & .993 & .910 \\
    BioLORD      & .993 & .993 & 1.000 & .861 \\
    MedCPT       & .792 & .951 & .757 & .868 \\
    med-gte      & .382 & .139 & .590 & .785 \\
    \bottomrule
  \end{tabular*}
\end{table}


\subsection{MIMIC-IV}
Table~\ref{tab:mimic} reports pooled results for the two audited concepts, and
Table~\ref{tab:mimic_concept} gives the concept-level and exploratory pooled
breakdown. On the audited concepts, dense-model UPFNR ranges from .413 to .547, HCWE from
.320 to .467, and observation-detection AUROC from .452 to .501. The proposition-only control obtains UPFNR .06 by predicting positive on
94\% of $C^+O^-$ cases, but its macro cell accuracy is only .65 versus .74 for
the patient-text model. Patient text therefore improves structured-state
discrimination without resolving documentation-sensitive errors. Encoder
rankings also differ from SilentEHR: Qwen has the lowest dense MIMIC UPFNR
despite one of the highest controlled UNR values. Results do not claim that
unavailable laboratories should be inferred; they show that LLM-labeled textual
absence is often mapped to clinical negativity.


\begin{table}[t]
  \centering
  \scriptsize
  \setlength{\tabcolsep}{2.2pt}
  \renewcommand{\arraystretch}{1.06}
  \caption{MIMIC-IV results on the audited concepts. Brackets show the
  subject-bootstrap 95\% CI for UPFNR.}
  \label{tab:mimic}
  \begin{tabular*}{\columnwidth}{@{\extracolsep{\fill}}lcccc@{}}
    \toprule
    \textbf{Encoder} & \textbf{UPFNR [CI]} &
    \textbf{HCWE} & \textbf{Gap} & \textbf{Det.} \\
    \midrule
    TF--IDF      & .653 [.49,.81] & .453 & .543 & .574 \\
    ClinicalBERT & .547 [.39,.69] & .467 & .207 & .452 \\
    Qwen3-Emb    & .440 [.29,.59] & .320 & .181 & .501 \\
    Jina v5      & .413 [.27,.56] & .320 & .158 & .480 \\
    BioLORD      & .533 [.38,.67] & .347 & .214 & .486 \\
    MedCPT       & .533 [.38,.68] & .427 & .251 & .459 \\
    med-gte      & .493 [.35,.66] & .347 & .166 & .474 \\
    \bottomrule
  \end{tabular*}
\end{table}



\begin{table*}[t]
  \centering
  \footnotesize
  \renewcommand{\arraystretch}{1.07}
  \caption{UPFNR by MIMIC-IV concept and pooled subset. Hyperkalemia is
  exploratory because only five structured-positive $O^-$ cases are available.}
  \label{tab:mimic_concept}
  \begin{tabular*}{\textwidth}{@{\extracolsep{\fill}}lrrrrrrrr@{}}
    \toprule
    \textbf{Subset} & \textbf{$n$} & \textbf{TF--IDF} &
    \textbf{ClinBERT} & \textbf{Qwen} & \textbf{Jina} &
    \textbf{BioLORD} & \textbf{MedCPT} & \textbf{med-gte} \\
    \midrule
    Anemia           & 61 & .69 & .56 & .44 & .44 & .48 & .52 & .49 \\
    Thrombocytopenia & 14 & .50 & .50 & .43 & .29 & .79 & .57 & .50 \\
    Audited pooled   & 75 & .653 & .547 & .440 & .413 & .533 & .533 & .493 \\
    \midrule
    Hyperkalemia     &  5 & .60 & 1.00 & .60 & .80 & 1.00 & 1.00 & .80 \\
    All three pooled & 80 & .650 & .575 & .450 & .438 & .562 & .562 & .512 \\
    \bottomrule
  \end{tabular*}
\end{table*}


Subject-bootstrap intervals exclude zero for all reported UPFNR estimates.
At the conservative label-precision setting $q=.60$, worst-case lower bounds
remain positive but approach zero for Qwen (.07) and Jina (.02), leaving their
noise-adjusted magnitudes uncertain. Because anemia contributes 61 of the 75
audited positive-$O^-$ cases and hyperkalemia has only five, we emphasize the
cross-model pattern and confidence-detection failure rather than small ranking
differences. These cells contain audited but disputed observation labels, not
gold textual silence.

\subsection{Human audit}
Table~\ref{tab:audit} summarizes agreement among the two annotators and
Qwen3-14B. Human annotators agree strongly ($\kappa=.80$). Agreement with Qwen3-14B is
moderate for seven statuses but stronger after collapsing to mentioned versus
unmentioned. Consensus binary accuracy is 89.0\% with $\kappa=.58$. Consensus $O^-$
precision and recall are both .65; per-annotator precision is .55. Reweighting
to the test distribution gives precision .60 [95\% CI .42,.78]. The main
disagreement is temporal, especially admission-only versus current evidence,
rather than whether the concept appears anywhere. All seven consensus structured-positive cases labeled $O^-$ by Qwen are
human-confirmed, but the subset is too small for precise estimation. The audit
therefore supports the coarse observation distinction while showing why MIMIC
effect sizes must remain qualified; consensus labels are available only for the
sampled rows.


\begin{table}[t]
  \centering
  \scriptsize
  \setlength{\tabcolsep}{2.2pt}
  \renewcommand{\arraystretch}{1.06}
  \caption{Human agreement with Qwen3-14B observation labels. Entries in
  the middle columns are accuracy/$\kappa$.}
  \label{tab:audit}
  \begin{tabular*}{\columnwidth}{@{\extracolsep{\fill}}lccc@{}}
    \toprule
    \textbf{Comparison} & \textbf{7-class} &
    \textbf{Binary} & \textbf{$O^-$ P/R/F1} \\
    \midrule
    H1 vs LLM   & 53.7/.39 & 88.3/.53 & .55/.65/.60 \\
    H2 vs LLM   & 54.7/.39 & 87.3/.50 & .55/.60/.58 \\
    Mean vs LLM & 54.2/.39 & 87.8/--- & --- \\
    Consensus   & 59.2/.44 & 89.0/.58 & .65/.65/.65 \\
    H1 vs H2    & 85.0/.80 & 99.0/.96 & --- \\
    \bottomrule
  \end{tabular*}
\end{table}


\subsection{Robustness and mitigation}
Alternative heads and regularization preserve high UNR for Qwen, Jina, and
BioLORD (.71--.99), although HCUNR can change sharply after calibration. This
supports a robust directional effect but not a calibration-invariant confidence
claim. Thresholds were fitted using known training examples only: the
confidence threshold was the fifth percentile of training maximum scores, and
the Mahalanobis threshold was the 95th percentile of training distances.
These choices target 95\% training coverage; held-out known coverage is
reported separately. Confidence abstention rejects only 3.5--39.6\% of
Unmentioned inputs for the neural encoders, yet accepted-set Unmentioned UNR
remains $\geq$.83 for Qwen3-Emb, Jina v5, BioLORD, and MedCPT
(Table~\ref{tab:abstention})---rejection does not preferentially remove
high-confidence CWE. Mahalanobis rejection is near zero. Unknown supervision
produces a clear dose-dependent reduction in Refuted-directed extrapolation,
with most of the improvement emerging between 10 and 50 semantic groups
(Table~\ref{tab:mitigation}); Qwen Unmentioned UNR falls from .944 to .058
while known accuracy stays above .98 at 50 groups.


\begin{table}[t]
  \centering
  \scriptsize
  \setlength{\tabcolsep}{2pt}
  \renewcommand{\arraystretch}{1.05}
  \caption{Confidence and Mahalanobis abstention on new-background transfer.
  Thresholds fit on known training cases only (5th-percentile confidence;
  95th-percentile Mahalanobis distance), targeting 95\% training coverage.
  Held-out known coverage and accepted-set Unmentioned UNR
  $P(\hat{y}{=}R \mid U_{\mathrm{unm}}, \mathrm{not\ rejected})$ are reported
  separately.}
  \label{tab:abstention}
  \begin{tabular*}{\columnwidth}{@{\extracolsep{\fill}}lcccc@{}}
    \toprule
    \textbf{Encoder} & \textbf{Known cov.} &
    \textbf{Unm.\ reject} & \textbf{Acc.\ UnmRR} &
    \textbf{Maha.\ reject} \\
    \midrule
    TF--IDF      & 1.00 & .98 & .00 & .00 \\
    ClinicalBERT & .89 & .79 & .73 & .00 \\
    Qwen3-Emb    & .96 & .28 & .97 & .00 \\
    Jina v5      & .85 & .40 & 1.00 & .00 \\
    BioLORD      & .95 & .04 & 1.00 & .00 \\
    MedCPT       & .84 & .13 & .83 & .00 \\
    med-gte      & .95 & .72 & .05 & .25 \\
    \bottomrule
  \end{tabular*}
\end{table}


\begin{table}[t]
  \centering
  \scriptsize
  \setlength{\tabcolsep}{2pt}
  \renewcommand{\arraystretch}{1.05}
  \caption{Unknown supervision on new-background transfer (10-seed mean).
  Base column is Unmentioned UNR under binary training.}
  \label{tab:mitigation}
  \begin{tabular*}{\columnwidth}{@{\extracolsep{\fill}}lcccc@{}}
    \toprule
    \textbf{Encoder} & \textbf{Base} &
    \textbf{10 grp} & \textbf{50 grp} & \textbf{100 grp} \\
    \midrule
    \multicolumn{5}{l}{\emph{UNR}} \\
    Qwen3-Emb & .944 & .410 & .058 & .026 \\
    Jina v5   & .910 & .539 & .224 & .112 \\
    BioLORD   & .993 & .672 & .294 & .135 \\
    MedCPT    & .792 & .433 & .282 & .201 \\
    med-gte   & .382 & .372 & .030 & .011 \\
    \midrule
    \multicolumn{5}{l}{\emph{Unknown recall}} \\
    Qwen3-Emb & --- & .517 & .888 & .923 \\
    Jina v5   & --- & .129 & .667 & .857 \\
    BioLORD   & --- & .162 & .587 & .785 \\
    MedCPT    & --- & .103 & .513 & .714 \\
    med-gte   & --- & .288 & .919 & .984 \\
    \midrule
    \multicolumn{5}{l}{\emph{Known accuracy}} \\
    Qwen3-Emb & --- & .983 & .982 & .983 \\
    Jina v5   & --- & .960 & .962 & .962 \\
    BioLORD   & --- & .988 & .989 & .986 \\
    MedCPT    & --- & .892 & .894 & .899 \\
    med-gte   & --- & .991 & .973 & .972 \\
    \bottomrule
  \end{tabular*}
\end{table}


Held-out-concept transfer is harder but remains effective for Qwen
(UNR .20; Unknown recall .75).

\section{Related Work}
\subsection{Clinical inference and assertion status}
General NLI benchmarks model entailment, contradiction, and neutral relations
\citep{bowman-etal-2015-large,williams-etal-2018-broad}. MedNLI transfers this
formulation to patient histories \citep{romanov-shivade-2018-lessons}, while
NLI4CT extends clinical inference to long trial reports and evidence selection
\citep{jullien-etal-2023-nli4ct}. These resources train and evaluate the neutral
class directly. Clinical assertion detection instead assigns mentioned entities
statuses such as present, possible, or absent
\citep{van-aken-etal-2021-assertion}; earlier work also extracts explicit
``evidence of absence'' from notes \citep{carter-etal-2016-evidence}. SilentEHR
addresses a different regime: the target proposition is not observed, and the
downstream classifier has been trained only on Supported and Refuted cases. The
question is therefore not whether a supervised Neutral label can be recognized,
but where a closed-world decision boundary sends non-observation.

\subsection{Clinical and biomedical representations}
Domain-specific pretraining has produced clinical contextual encoders such as
ClinicalBERT \citep{alsentzer-etal-2019-clinicalbert}, ontology-grounded
biomedical sentence representations such as BioLORD
\citep{remy-etal-2022-biolord}, and search-log-trained biomedical retrievers such
as MedCPT \citep{jin-etal-2023-medcpt}. Recent general embedding models use
large-scale contrastive training, instruction conditioning, and distillation,
including Qwen3-Embedding \citep{zhang-etal-2025-qwen3embedding} and Jina v5
\citep{akram-etal-2026-jinav5}. Their evaluations emphasize semantic
similarity, retrieval, clustering, or classification. Our analysis is
complementary: encoder weights remain frozen, and we audit how a supervised
binary head extrapolates to an epistemic state excluded from training. Strong
benchmark embeddings can therefore still exhibit unsafe closed-world behavior.

\subsection{Missingness, selective prediction, and OOD detection}
Missingness in structured EHR time series can itself be informative and has
been modeled explicitly, for example with GRU-D \citep{che-etal-2018-grud}.
Textual non-observation is related but distinct: a concept can be absent from a
note because it was not assessed, not relevant to that document, historical, or
simply omitted. Selective prediction and OOD detection commonly use maximum
class probability \citep{hendrycks-gimpel-2017-baseline}, calibrated confidence
\citep{guo-etal-2017-calibration}, learned reject options
\citep{geifman-el-yaniv-2019-selectivenet}, or class-conditional Mahalanobis
distance \citep{lee-etal-2018-mahalanobis}. SilentEHR tests a harder boundary
case for these methods: Unmentioned profiles are linguistically and clinically
in-domain, but their evidence-coverage state is outside the binary training
support. Our findings show why low confidence or global distance need not track
the semantic direction induced by the downstream head.

\section{Discussion}
CWE is primarily a directional extrapolation failure, not a universal geometric
collapse. Its severity depends on encoder, epistemic state, concept family, and
evidence realization. MIMIC-IV shows a related real-note consequence, but
encoder rankings differ from SilentEHR and observation labels remain imperfect.
Accordingly, MIMIC $C^+O^-$ cells should be interpreted as structured
positives with disputed textual silence, not as a gold set of undocumented
facts. The controlled benchmark supplies internal validity, whereas MIMIC-IV
tests ecological relevance under noisier observation labels. Explicit Unknown
supervision is more effective than generic confidence or distance rejection,
suggesting that coverage state must be learned rather than recovered from a
binary model after deployment.

\section{Conclusion}
Frozen binary verifiers often map non-observation to Refuted, but the effect
varies across encoders, concepts, and epistemic states. SilentEHR separates
direction from confidence and geometry, while MIMIC-IV demonstrates a related
documentation-mismatch risk in real notes. Clinical verification systems should
represent and supervise non-observation explicitly.

\section*{Ethics Statement}
MIMIC-IV and MIMIC-IV-Note are used under their data-use agreements; all notes
are de-identified. Human annotations were collected only for label validation
by two graduate-student annotators working independently under a blind protocol
(Appendix~\ref{app:annotation_protocol}); no third adjudicator was involved.
Disagreements were excluded from consensus gold rather than resolved by a senior
reviewer.

\section*{Limitations}
SilentEHR uses synthetic backgrounds and 16 concepts. MIMIC-IV covers three
near-discharge laboratory concepts, while the human audit covers two.
Qwen3-14B observation labels remain imperfect; 45 original audit disagreements
were excluded rather than adjudicated by a third annotator, and the
human-confirmed positive-$O^-$ subset is small.
Hyperkalemia estimates are unstable because only five positive $O^-$ cases are
available. Finally, Unknown supervision requires labeled examples or semantic
groups that may be unavailable in deployment.

\bibliography{silent_ehr_refs}
\bibliographystyle{acl_natbib}

\appendix
\section{Evaluation Protocols}
\label{app:protocols}
Table~\ref{tab:protocols} summarizes the four SilentEHR split protocols.

\begin{table}[!htbp]
  \centering
  \footnotesize
  \renewcommand{\arraystretch}{1.08}
  \caption{SilentEHR evaluation protocols.}
  \label{tab:protocols}
  \begin{tabular*}{\columnwidth}{@{\extracolsep{\fill}}lcc@{}}
    \toprule
    \textbf{Protocol} & \textbf{Held out} & \textbf{U at test?} \\
    \midrule
    New background & 12/80 backgrounds & Yes \\
    IID & background$\times$concept groups & Yes \\
    New template & assertion templates & No \\
    New concept & 4/16 concepts & Yes \\
    \bottomrule
  \end{tabular*}
\end{table}

\section{MIMIC State--Observation Cell Design}
\label{app:mimic_controls_detail}
The four-cell design used in the structured-state MIMIC analysis
(Appendix~\ref{app:mimic_controls}) is defined as follows. Training cells:
$C^+O^+$, $C^-O^+$.
Evaluation adds $C^+O^-$ and $C^-O^-$ where structured laboratories indicate
abnormality/normality but Qwen3-14B marks the concept \emph{unmentioned} in the
discharge note.
Historical-only, uncertain, admission-only, and temporal-conflict Qwen3-14B
labels are excluded from this binary protocol (but map to Unknown in the
separate direct S/R/U bridge; Appendix~\ref{app:real_sru}).

\section{Per-Concept Unmentioned UNR}
\label{app:concept_unr}

Twelve concepts appear in the new-background training split; four held-out
concepts are evaluated under the new-concept protocol (Table~\ref{tab:protocols}).
Qwen3-Embedding, Jina v5, and BioLORD predict Refuted for Unmentioned inputs in
all 12 training concepts, whereas med-gte exceeds 0.5 in only four. The med-gte
contrast is concentrated in allergy and medication concepts, while disease
concepts remain strongly negative (Tables~\ref{tab:concept_unr} and
\ref{tab:concept_group}). Its allergy advantage disappears under the
indirect-evidence tier, indicating evidence-realization sensitivity rather than
a stable abstention policy.

\begin{table*}[t]
  \centering
  \scriptsize
  \renewcommand{\arraystretch}{1.04}
  \caption{Clean-tier Unmentioned-to-Refuted rate by concept.}
  \label{tab:concept_unr}
  \begin{tabular*}{\textwidth}{@{\extracolsep{\fill}}lrrrrrrr@{}}
    \toprule
    \textbf{Concept} & \textbf{TF--IDF} & \textbf{ClinBERT} &
    \textbf{Qwen} & \textbf{Jina} & \textbf{BioLORD} &
    \textbf{MedCPT} & \textbf{med-gte} \\
    \midrule
    Chest pain & .92 & .75 & .83 & 1.00 & 1.00 & .17 & .50 \\
    Fever & .92 & .33 & 1.00 & 1.00 & 1.00 & .08 & .08 \\
    Penicillin allergy & .00 & .17 & 1.00 & 1.00 & 1.00 & 1.00 & .00 \\
    Active bleeding & .92 & .92 & 1.00 & 1.00 & 1.00 & 1.00 & .58 \\
    Anticoagulation & .92 & .92 & .92 & .75 & 1.00 & 1.00 & .08 \\
    Renal dysfunction & .92 & .92 & .92 & 1.00 & .92 & 1.00 & .83 \\
    Nausea & .92 & .17 & 1.00 & 1.00 & 1.00 & .75 & 1.00 \\
    Sulfonamide allergy & .00 & .83 & .92 & 1.00 & 1.00 & 1.00 & .00 \\
    Latex allergy & .00 & .17 & 1.00 & 1.00 & 1.00 & 1.00 & .00 \\
    Insulin therapy & .92 & .17 & .92 & 1.00 & 1.00 & 1.00 & .00 \\
    Antiplatelet therapy & .92 & .92 & .83 & .58 & 1.00 & 1.00 & .50 \\
    Cognitive impairment & .92 & .92 & 1.00 & .58 & 1.00 & .50 & 1.00 \\
    \midrule
    Concepts $>.5$ & 9 & 7 & 12 & 12 & 12 & 9 & 4 \\
    \bottomrule
  \end{tabular*}
\end{table*}

\begin{table}[!htbp]
  \centering
  \footnotesize
  \renewcommand{\arraystretch}{1.08}
  \caption{Mean Unmentioned-to-Refuted rate by concept family.}
  \label{tab:concept_group}
  \begin{tabular*}{\columnwidth}{@{\extracolsep{\fill}}lrrrr@{}}
    \toprule
    \textbf{Family} & \textbf{Qwen} & \textbf{BioLORD} &
    \textbf{med-gte} & \textbf{gte-ind.} \\
    \midrule
    Symptom    & .94 & 1.00 & .53 & 1.00 \\
    Allergy    & .97 & 1.00 & .00 & 1.00 \\
    Medication & .89 & 1.00 & .19 & .39 \\
    Disease    & .96 & .96 & .92 & 1.00 \\
    Safety     & 1.00 & 1.00 & .58 & .25 \\
    \bottomrule
  \end{tabular*}
\end{table}

\FloatBarrier

\section{Unknown-State Decision-Axis Logits}
\label{app:unknown_logit}

Positive logits favor Supported and negative logits favor Refuted. Qwen3-Embedding
and BioLORD place both Unmentioned and Not-assessed strongly on the Refuted side.
Jina v5 instead separates explicit uncertainty from silence: its mean logit is
positive for Uncertain but strongly negative for Unmentioned. These results show
that the three Unknown states are represented differently even when the binary
head maps several of them to the same class (Table~\ref{tab:unknown_logit}).

\begin{table}[!htbp]
  \centering
  \scriptsize
  \renewcommand{\arraystretch}{1.06}
  \caption{Mean decision-axis logit; positive values favor Supported.}
  \label{tab:unknown_logit}
  \begin{tabular*}{\columnwidth}{@{\extracolsep{\fill}}lrrrrr@{}}
    \toprule
    \textbf{Encoder} & \textbf{S} & \textbf{R} &
    \textbf{U} & \textbf{NA} & \textbf{Unc} \\
    \midrule
    TF--IDF      & 9.20 & $-8.74$ & $-1.05$ & .37 & $-1.04$ \\
    ClinicalBERT & 7.55 & $-6.50$ & $-1.02$ & .65 & $-.25$ \\
    Qwen3-Emb    & 11.31 & $-9.72$ & $-7.85$ & $-9.27$ & $-2.65$ \\
    Jina v5      & 5.53 & $-7.83$ & $-8.61$ & $-1.29$ & 5.04 \\
    BioLORD      & 8.99 & $-10.34$ & $-16.50$ & $-3.35$ & $-2.85$ \\
    MedCPT       & 3.60 & $-6.99$ & $-11.68$ & 3.07 & $-.32$ \\
    med-gte      & 7.59 & $-8.00$ & 4.64 & $-.16$ & $-.16$ \\
    \bottomrule
  \end{tabular*}
\end{table}

\FloatBarrier

\section{Split Stability}
\label{app:split_stability}

The primary results are not driven by one favorable partition. Across ten
background splits, Qwen3-Embedding, Jina v5, BioLORD, MedCPT, and TF--IDF exceed an
Unmentioned-to-Refuted rate of 0.5 in every split, whereas med-gte remains below
0.5 throughout and ClinicalBERT exceeds 0.5 in eight of ten splits. Four held-out-concept folds reproduce the same encoder-level
ordering. Known-state accuracy remains high, so the directional differences
cannot be attributed to unstable recognition of explicit evidence
(Table~\ref{tab:split_stability}).

\begin{table}[!htbp]
  \centering
  \scriptsize
  \renewcommand{\arraystretch}{1.06}
  \caption{Split stability. Values are mean $\pm$ SD; parentheses give
  splits with U$\to$R $>.5$.}
  \label{tab:split_stability}
  \begin{tabular*}{\columnwidth}{@{\extracolsep{\fill}}lcccc@{}}
    \toprule
    & \multicolumn{2}{c}{\textbf{Background}} &
    \multicolumn{2}{c}{\textbf{Concept}} \\
    \cmidrule(lr){2-3}\cmidrule(lr){4-5}
    \textbf{Encoder} & \textbf{Known} & \textbf{U$\to$R} &
    \textbf{Known} & \textbf{U$\to$R} \\
    \midrule
    TF--IDF & 1.00$\pm$.00 & .71$\pm$.04 (10) &
              1.00$\pm$.00 & .88$\pm$.11 (4) \\
    Qwen    & .93$\pm$.04 & .90$\pm$.06 (10) &
              .95$\pm$.04 & .98$\pm$.01 (4) \\
    Jina    & .93$\pm$.02 & .85$\pm$.04 (10) &
              .85$\pm$.06 & .88$\pm$.02 (4) \\
    BioLORD & .95$\pm$.02 & .97$\pm$.03 (10) &
              .89$\pm$.06 & .99$\pm$.01 (4) \\
    med-gte & .93$\pm$.04 & .34$\pm$.08 (0) &
              .84$\pm$.19 & .24$\pm$.16 (0) \\
    MedCPT  & .85$\pm$.03 & .75$\pm$.05 (10) &
              .71$\pm$.08 & .79$\pm$.08 (4) \\
    ClinicalBERT & .93$\pm$.04 & .58$\pm$.11 (8) &
              .99$\pm$.01 & .67$\pm$.18 (4) \\
    \bottomrule
  \end{tabular*}
\end{table}

\FloatBarrier

\section{MIMIC Negative Controls}
\label{app:mimic_controls}
Table~\ref{tab:mimic_controls} shows that low UPFNR can be achieved by a
degenerate positive bias; the four-cell recalls expose this failure.

\begin{table}[!htbp]
  \centering
  \scriptsize
  \renewcommand{\arraystretch}{1.05}
  \caption{MIMIC negative controls using Qwen3-Embedding.}
  \label{tab:mimic_controls}
  \begin{tabular*}{\columnwidth}{@{\extracolsep{\fill}}lrrrr@{}}
    \toprule
    \multicolumn{5}{l}{\emph{Overall metrics}} \\
    \textbf{Model} & \textbf{Macro} & \textbf{UPFNR} &
    \textbf{HCWE} & \textbf{Det.} \\
    \midrule
    Patient + prop. & .74 & .43 & .38 & .48 \\
    Proposition only & .65 & .06 & .06 & .39 \\
    Concept prior & .65 & .06 & .06 & .45 \\
    \addlinespace[3pt]
    \multicolumn{5}{l}{\emph{Cell recall}} \\
    \textbf{Model} & \textbf{$C^+O^+$} & \textbf{$C^-O^+$} &
    \textbf{$C^+O^-$} & \textbf{$C^-O^-$} \\
    \midrule
    Patient + prop. & .73 & .84 & .57 & .82 \\
    Proposition only & .97 & .33 & .94 & .35 \\
    Concept prior & .97 & .33 & .94 & .35 \\
    \bottomrule
  \end{tabular*}
\end{table}

\section{Human-Audit Subset Sensitivity}
\label{app:audit_sensitivity}

We re-evaluate the MIMIC classifier on 300 dual-annotated admission--concept
cases (anemia and thrombocytopenia only) as an exploratory sensitivity analysis, reporting UPFNR ($=1-\text{Recall}(C^+)$) on $C^+$ audit rows under
five operational definitions of textual $O^-$: Human1, Human2, consensus,
union (either annotator marks \emph{unmentioned}), and intersection (both
annotators agree on \emph{unmentioned}).
Wilson 95\% intervals are reported for binomial proportions.
The Qwen3-14B $O^-$ $C^+$ subset ($n=7$) coincides with the
Qwen3-14B--annotator-agreement subset on this audit sample. Table~\ref{tab:audit_sensitivity} reports the
resulting estimates and Wilson intervals.

\begin{table}[!htbp]
  \centering
  \scriptsize
  \renewcommand{\arraystretch}{1.06}
  \caption{Audit-subset UPFNR [Wilson 95\% CI]. H1/consensus/intersection use
  $n=18$; H2/union use $n=19$; Qwen3-14B $O^-$ uses $n=7$.}
  \label{tab:audit_sensitivity}
  \begin{tabular*}{\columnwidth}{@{\extracolsep{\fill}}lccc@{}}
    \toprule
    \textbf{Encoder} & \textbf{H1/Cons./Int.} &
    \textbf{H2/Union} & \textbf{Qwen3-14B} \\
    \midrule
    TF--IDF & .44 [.25,.66] & .47 [.27,.68] & 1.00 [.65,1.00] \\
    Qwen    & .17 [.06,.39] & .16 [.06,.38] & .29 [.08,.64] \\
    Jina    & .28 [.12,.51] & .26 [.12,.49] & .57 [.25,.84] \\
    BioLORD & .33 [.16,.56] & .32 [.15,.54] & .71 [.36,.92] \\
    \bottomrule
  \end{tabular*}
\end{table}

\section{Audit-Calibrated Label-Noise Sensitivity}
\label{app:label_noise}

Human audit estimates that 55--65\% of Qwen3-14B \emph{unmentioned}
labels are confirmed (Appendix Table~\ref{tab:audit_ominus_precision}).
We report two complementary analyses on the audited-concept
$C^+O^-$ cohort ($N=75$), regenerated from the same per-concept classifier
protocol as Table~\ref{tab:mimic} (code asserts
$|\text{observed UPFNR} - \text{main table}| < 10^{-9}$).
Contamination assumptions: $q=0.55$ (per-annotator conservative),
$q=0.60$ (overall cohort-weighted $O^-$ precision, applied conservatively to
$C^+O^-$), and $q=0.65$ (consensus audit-sample optimistic).
Because the $C^+$ audit strata contain only seven Qwen3-14B-$O^-$ cases (all
human-confirmed), we avoid treating their observed precision of 1.00 as a
stable $C^+$-specific estimate and use the overall cohort-weighted value of
0.60 as the primary conservative setting.

\subsection{Missing-at-random scenario}
In each of 1{,}000 Monte Carlo replicates, we independently retain a fraction
$p_{\text{keep}}=q$ of Qwen-labeled $C^+O^-$ cases (concept-stratified) and
exclude the rest as putative false $O^-$ labels.
This assumes false labels are deleted independently of model prediction; median
UPFNR closely tracks the observed value in Table~\ref{tab:mimic}
(Table~\ref{tab:label_noise_mar}).

\subsection{Worst-case partial-identification bounds}
Let $N=75$ and $E$ be model-predicted-negative count on Qwen-labeled $C^+O^-$.
Assuming at least fraction $q$ of labels are true $O^-$, the true UPFNR lies in
\[
  \Bigl[\max\!\Bigl(0,\frac{E-(1-q)N}{qN}\Bigr),\;
        \min\!\Bigl(1,\frac{E}{qN}\Bigr)\Bigr],
\]
allowing false-$O^-$ labels to concentrate on model-error or model-correct
cases. Tables~\ref{tab:label_noise_bounds} and \ref{tab:label_noise_mar}
report the worst-case and missing-at-random results.

\begin{table}[!htbp]
  \centering
  \scriptsize
  \renewcommand{\arraystretch}{1.05}
  \caption{Worst-case UPFNR bounds; ranges are [LB, UB].}
  \label{tab:label_noise_bounds}
  \begin{tabular*}{\columnwidth}{@{\extracolsep{\fill}}lcccc@{}}
    \toprule
    \textbf{Encoder} & \textbf{Obs.} & \textbf{$q=.55$} &
    \textbf{$q=.60$} & \textbf{$q=.65$} \\
    \midrule
    Qwen3-Emb    & .440 & [.00,.80] & [.07,.73] & [.14,.68] \\
    Jina v5      & .413 & [.00,.75] & [.02,.69] & [.10,.64] \\
    BioLORD      & .533 & [.15,.97] & [.22,.89] & [.28,.82] \\
    ClinicalBERT & .547 & [.18,.99] & [.24,.91] & [.30,.84] \\
    MedCPT       & .533 & [.15,.97] & [.22,.89] & [.28,.82] \\
    med-gte      & .493 & [.08,.90] & [.16,.82] & [.22,.76] \\
    TF--IDF      & .653 & [.37,1] & [.42,1] & [.47,1] \\
    \bottomrule
  \end{tabular*}
\end{table}

\begin{table}[!htbp]
  \centering
  \footnotesize
  \renewcommand{\arraystretch}{1.08}
  \caption{MAR label-noise sensitivity at $p_{\mathrm{keep}}=.60$
  (1,000 replicates).}
  \label{tab:label_noise_mar}
  \begin{tabular*}{\columnwidth}{@{\extracolsep{\fill}}lccc@{}}
    \toprule
    \textbf{Encoder} & \textbf{Median} & \textbf{95\% range} &
    \textbf{HCWE} \\
    \midrule
    Qwen3-Emb    & .435 & [.341, .529] & .349 \\
    Jina v5      & .415 & [.326, .512] & .360 \\
    BioLORD      & .533 & [.435, .625] & .347 \\
    ClinicalBERT & .546 & [.454, .644] & .468 \\
    MedCPT       & .532 & [.442, .623] & .429 \\
    med-gte      & .500 & [.395, .591] & .347 \\
    TF--IDF      & .649 & [.565, .744] & .450 \\
    \bottomrule
  \end{tabular*}
\end{table}

\section{Stratified Human Precision for Unmentioned Labels}
\label{app:audit_ominus_precision}

Table~\ref{tab:audit_ominus_precision} decomposes $P(\text{human }O^- \mid
\text{Qwen }O^-)$ on the 300-case audit sample and reweights stratum
estimates by the audited-concept MIMIC test distribution of Qwen3-14B-$O^-$
cells ($N_{\text{cohort}}=805$).

\begin{table}[!htbp]
  \centering
  \footnotesize
  \renewcommand{\arraystretch}{1.08}
  \caption{Consensus precision of Qwen-\emph{unmentioned} labels by audit
  stratum. Weights follow the audited-concept test distribution.}
  \label{tab:audit_ominus_precision}
  \begin{tabular*}{\columnwidth}{@{\extracolsep{\fill}}lrrr@{}}
    \toprule
    \textbf{Stratum} & \textbf{$n$} & \textbf{Prec.} & \textbf{Weight} \\
    \midrule
    Anemia $C^-$           & 4  & .50 & .29 \\
    Anemia $C^+$           & 3  & 1.00 & .08 \\
    Thrombocytopenia $C^-$ & 29 & .59 & .62 \\
    Thrombocytopenia $C^+$ & 4  & 1.00 & .02 \\
    \midrule
    Pooled audit           & 40 & .65 & --- \\
    Cohort weighted        & --- & .60 & --- \\
    Per-annotator pooled   & 47 & .55 & --- \\
    \bottomrule
  \end{tabular*}
  \vspace{2pt}
  {\scriptsize Cohort-weighted 95\% bootstrap CI: [.42, .78].}
\end{table}

\FloatBarrier

\section{SilentEHR Input Ablation}
\label{app:input_ablation}

Table~\ref{tab:input_ablation} isolates whether CWE arises from patient
embeddings, proposition embeddings, or their pairing. Proposition-only features
collapse to chance on Known cases (0.50) and zero Unmentioned UNR because the
head cannot distinguish patients. Patient-only features preserve high UNR for
Qwen3-Embedding, Jina v5, and BioLORD, showing that patient geometry alone
drives much of the collapse. Shuffling proposition assignments at test time
preserves UNR but degrades Known accuracy, confirming that correct pairing is
required for evidence recognition yet unnecessary for silence$\to$Refuted
mapping.

\begin{table}[!htbp]
  \centering
  \scriptsize
  \renewcommand{\arraystretch}{1.06}
  \caption{Input ablation on new-background protocol (Known acc /
  Unmentioned UNR).}
  \label{tab:input_ablation}
  \begin{tabular*}{\columnwidth}{@{\extracolsep{\fill}}lcccccc@{}}
    \toprule
    & \multicolumn{2}{c}{\textbf{Full}} &
    \multicolumn{2}{c}{\textbf{Patient}} &
    \multicolumn{2}{c}{\textbf{Prop / Shuf}} \\
    \cmidrule(lr){2-3}\cmidrule(lr){4-5}\cmidrule(lr){6-7}
    \textbf{Encoder} & \textbf{K} & \textbf{U} &
    \textbf{K} & \textbf{U} & \textbf{Prop K} & \textbf{Shuf K} \\
    \midrule
    Qwen3-Emb & .98 & .94 & .96 & .92 & .50 & .79 \\
    Jina v5   & .96 & .91 & .96 & 1.00 & .50 & .69 \\
    BioLORD   & .99 & .99 & .98 & 1.00 & .50 & .64 \\
    med-gte   & 1.00 & .38 & .97 & .08 & .50 & .79 \\
    TF--IDF   & 1.00 & .69 & 1.00 & .92 & .50 & 1.00 \\
    \bottomrule
  \end{tabular*}
\end{table}

\section{Pair-Feature Composition Ablation}
\label{app:pair_ablation}

Table~\ref{tab:pair_ablation} varies the pair representation fed to the
Supported/Refuted head. For Qwen3-Embedding and Jina v5, adding
$|e_x-e_h|$ and $e_x\odot e_h$ \emph{increases} Unmentioned UNR relative to
$e_x$ alone, while Known accuracy remains high. CWE is therefore not explained
by the linear head alone: interaction features amplify collapse once patient
context is available.

\begin{table}[!htbp]
  \centering
  \scriptsize
  \renewcommand{\arraystretch}{1.06}
  \caption{Pair-feature ablation (Known / Unmentioned UNR / HCUNR@0.9 /
  Unmentioned-detection AUROC). Full = $[e_x;e_h;|e_x-e_h|;e_x\odot e_h]$.}
  \label{tab:pair_ablation}
  \begin{tabular*}{\columnwidth}{@{\extracolsep{\fill}}llcccc@{}}
    \toprule
    \textbf{Encoder} & \textbf{Features} & \textbf{K} & \textbf{U} &
    \textbf{HC} & \textbf{Det.} \\
    \midrule
    Qwen3-Emb & $e_x$ & .96 & .92 & .75 & .96 \\
                & $[e_x;e_h]$ & .96 & .72 & .42 & .91 \\
                & Full & .98 & .94 & .88 & .66 \\
    Jina v5   & $e_x$ & .96 & 1.00 & .42 & .92 \\
                & $[e_x;e_h]$ & .96 & .68 & .54 & .78 \\
                & Full & .96 & .91 & .81 & .50 \\
    BioLORD   & $e_x$ & .98 & 1.00 & 1.00 & .70 \\
                & $[e_x;e_h]$ & .97 & .92 & .88 & .60 \\
                & Full & .99 & .99 & .99 & .18 \\
    med-gte   & $e_x$ & .97 & .08 & .08 & .94 \\
                & $[e_x;e_h]$ & .96 & .10 & .07 & .75 \\
                & Full & 1.00 & .38 & .16 & .75 \\
    \bottomrule
  \end{tabular*}
\end{table}

\section{Mechanism Probes on Clean SilentEHR}
\label{app:mechanism_clean}

Local Refuted-neighborhood composition (Local-R@10) correlates with UNR for
high-collapse encoders: BioLORD reaches Local-R@10$=.72$ with UNR$=.99$.
Linear Unmentioned-vs-known probes achieve AUROC$\approx$1.0 on clean
background-only Unmentioned cases, indicating strong separability of silence
from explicit evidence in embedding space (Table~\ref{tab:mechanism_clean}).

\begin{table}[!htbp]
  \centering
  \scriptsize
  \renewcommand{\arraystretch}{1.06}
  \caption{Mechanism probes on clean new-background split.}
  \label{tab:mechanism_clean}
  \begin{tabular*}{\columnwidth}{@{\extracolsep{\fill}}lrrr@{}}
    \toprule
    \textbf{Encoder} & \textbf{UNR} & \textbf{Local-R@10} &
    \textbf{U/K AUROC} \\
    \midrule
    Qwen3-Emb    & .94 & .65 & 1.00 \\
    Jina v5      & .91 & .57 & .98 \\
    BioLORD      & .99 & .72 & .99 \\
    med-gte      & .38 & .52 & 1.00 \\
    TF--IDF      & .69 & .51 & 1.00 \\
    ClinicalBERT & .60 & .52 & 1.00 \\
    MedCPT       & .79 & .53 & 1.00 \\
    \bottomrule
  \end{tabular*}
\end{table}

\section{Mechanism Probe Robustness Across Tiers}
\label{app:mechanism_robust}

Table~\ref{tab:mechanism_robust} repeats U/K AUROC and Local-R@10 on
held-concept, indirect-evidence, and length-matched tiers. Separability
remains high ($\geq$.97 AUROC) under indirect and length-matched controls;
Local-R@10 stays above chance ($>.50$) for dense encoders in all tiers.

\begin{table}[!htbp]
  \centering
  \scriptsize
  \renewcommand{\arraystretch}{1.06}
  \caption{Mechanism probe robustness (Qwen3-Embedding shown; others similar).}
  \label{tab:mechanism_robust}
  \begin{tabular*}{\columnwidth}{@{\extracolsep{\fill}}lccc@{}}
    \toprule
    \textbf{Tier} & \textbf{Qwen U/K} & \textbf{Qwen LR@10} &
    \textbf{BioLORD LR@10} \\
    \midrule
    Held concept      & 1.00 & .58 & .64 \\
    Indirect evidence & 1.00 & .68 & .62 \\
    Length matched    & 1.00 & .62 & .72 \\
    \bottomrule
  \end{tabular*}
\end{table}

\section{HCUNR Threshold Sensitivity}
\label{app:hcunr_sensitivity}

Table~\ref{tab:hcunr_sensitivity} reports high-confidence UNR at
$\tau\in\{0.8,0.9,0.95\}$. Rankings are stable: BioLORD and Qwen3-Embedding
remain above 0.83 at $\tau=0.95$; med-gte stays below 0.10.

\begin{table}[!htbp]
  \centering
  \scriptsize
  \renewcommand{\arraystretch}{1.06}
  \caption{Unmentioned HCUNR at alternative confidence thresholds.}
  \label{tab:hcunr_sensitivity}
  \begin{tabular*}{\columnwidth}{@{\extracolsep{\fill}}lcccc@{}}
    \toprule
    \textbf{Encoder} & \textbf{UNR} & \textbf{@.8} & \textbf{@.9} &
    \textbf{@.95} \\
    \midrule
    Qwen3-Emb    & .94 & .92 & .88 & .83 \\
    Jina v5      & .91 & .83 & .81 & .76 \\
    BioLORD      & .99 & .99 & .99 & .99 \\
    med-gte      & .38 & .22 & .16 & .10 \\
    TF--IDF      & .69 & .62 & .46 & .00 \\
    ClinicalBERT & .60 & .38 & .27 & .19 \\
    MedCPT       & .79 & .74 & .72 & .72 \\
    \bottomrule
  \end{tabular*}
\end{table}

\section{MIMIC Cell Recalls and Temporal Observability}
\label{app:mimic_cells}

Table~\ref{tab:mimic_cell_recalls} reports per-cell recall
$R=P(\hat C=C\mid C,O)$ on the matched MIMIC cohort. All encoders achieve
$R^{-, -}>0.76$ on $C^-O^-$ cases, showing that negative structured states
are recognized when documentation is absent; UPFNR on $C^+O^-$ therefore reflects
asymmetric extrapolation rather than blanket negative bias.

\begin{table}[!htbp]
  \centering
  \scriptsize
  \renewcommand{\arraystretch}{1.06}
  \caption{MIMIC cell recalls ($n=1122$ for $C^-O^-$; $n=80$ for $C^+O^-$).}
  \label{tab:mimic_cell_recalls}
  \begin{tabular*}{\columnwidth}{@{\extracolsep{\fill}}lrrr@{}}
    \toprule
    \textbf{Encoder} & \textbf{$R^{+,+}$} & \textbf{$R^{+,-}$} &
    \textbf{$R^{-,-}$} \\
    \midrule
    Qwen3-Emb    & .73 & .58 & .82 \\
    Jina v5      & .71 & .51 & .83 \\
    BioLORD      & .66 & .46 & .76 \\
    ClinicalBERT & .65 & .40 & .81 \\
    MedCPT       & .69 & .53 & .82 \\
    med-gte      & .67 & .54 & .82 \\
    TF--IDF      & .82 & .19 & .85 \\
    \bottomrule
  \end{tabular*}
\end{table}

Discharge-note \emph{charttime} and \emph{storetime} are not exported in our
MIMIC struct cohort; we therefore use \emph{dischtime} as a conservative proxy
for note availability. Labs are drawn from the final 24\,h before discharge, so
$t_{\mathrm{lab}}\le t_{\mathrm{discharge}}$ holds for 100\% (80/80) of
$C^+O^-$ test cases. Under this conservative subset,
$R^{+,-}=0.58$ (Qwen3-Embedding), unchanged from the full $C^+O^-$ subset.

\section{Real Clinical Supported/Refuted/Unknown Validation}
\label{app:real_sru}

We construct a human-gold S/R/U bridge from the 300-case MIMIC audit cohort
with dual annotators (Human1 + Human2). Seven-class observation labels map to
S/R/U as follows: \emph{observed\_positive}$\to$Supported,
\emph{observed\_negative}$\to$Refuted, and
\emph{unmentioned}/\emph{uncertain}/\emph{historical\_only}/\emph{admission\_only}/\emph{temporal\_conflict}$\to$Unknown.
Dual-annotator consensus holds on 255/300 original audit pairs
($\kappa\approx$.80); 45 disagreements are excluded (no third adjudicator).
A 50-pair non-laboratory expansion (strict \emph{unmentioned} propositions on
the same notes) adds 50/50 consensus rows, yielding 305 consensus pairs total.
Table~\ref{tab:data_flow} summarizes how this cohort relates to the 493-pair
silver set and the Qwen3-14B labeling pool.

Our primary estimate uses grouped subject-level 5-fold cross-validation: train
Supported+Refuted only on four folds and report pooled out-of-fold metrics on
all 305 consensus rows ($n_{\mathrm{Unknown}}=159$; 83 strict-\emph{unmentioned},
62 \emph{admission-only}; Wilson 95\% intervals in
Table~\ref{tab:real_sru_cv}). A single 70/30 subject split on the original
255 lab-audit consensus rows ($n_{\mathrm{test}}=77$, 34 Unknown; seed~0) is
retained as a sensitivity check (Table~\ref{tab:real_sru_split}) because split
parity was verified on that subset
(\texttt{outputs/real\_clinical\_sru\_split\_parity.json}); all models pass
direction sanity ($\bar{p}_{\mathrm{ref}}|R > \bar{p}_{\mathrm{ref}}|S$ on
known training rows) with explicit \texttt{clf.classes\_} indexing.

On pooled OOF consensus gold, frozen verifiers do \emph{not} systematically
collapse silence to Refuted: Qwen3-Emb UNR$=.37$ [.30,.45], Jina v5
UNR$=.36$ [.29,.44], BioLORD UNR$=.35$ [.28,.43]---far below SilentEHR
Unmentioned UNR ($\geq$.91 for the same encoders). Strict-\emph{unmentioned}
Unknown rows show higher UNR (.47--.52) than \emph{admission-only} rows
(.13--.26), consistent with controlled absence being the strongest CWE setting.
This supports treating CWE as strongest under controlled omission while real
discharge-note silence is often mapped to Unknown rather than Refuted under
human-consensus S/R/U labels. The 493-pair silver set
(\texttt{real\_clinical\_sru\_silver.jsonl}) adds Qwen3-14B-mapped pairs on the
same notes for sensitivity only (Appendix~\ref{app:data_flow}).

\begin{table}[!htbp]
  \centering
  \scriptsize
  \renewcommand{\arraystretch}{1.06}
  \caption{Real clinical S/R/U (human consensus; grouped 5-fold CV; train S/R
  only). Pooled out-of-fold Wilson 95\% intervals.}
  \label{tab:real_sru_cv}
  \begin{tabular*}{\columnwidth}{@{\extracolsep{\fill}}lrrrrr@{}}
    \toprule
    \textbf{Encoder} & \textbf{Known} & \textbf{UNR} &
    \textbf{HC@.9} & \textbf{Adm.} & \textbf{Strict unm.} \\
    \midrule
    Qwen3-Emb    & .67 [.59,.74] & .37 [.30,.45] &
    .34 [.24,.44] & .26 [.17,.38] & .47 [.37,.58] \\
    Jina v5      & .70 [.62,.77] & .36 [.29,.44] &
    .36 [.27,.47] & .23 [.14,.34] & .49 [.39,.60] \\
    BioLORD      & .67 [.59,.74] & .35 [.28,.43] &
    .33 [.23,.43] & .15 [.08,.25] & .52 [.41,.62] \\
    TF--IDF      & .63 [.55,.70] & .33 [.26,.40] &
    .29 [.20,.40] & .13 [.07,.23] & .47 [.37,.58] \\
    \bottomrule
  \end{tabular*}
\end{table}

\begin{table}[!htbp]
  \centering
  \scriptsize
  \renewcommand{\arraystretch}{1.06}
  \caption{Single 70/30 subject split sensitivity on 255 lab-audit consensus
  rows (train S/R only; seed~0).}
  \label{tab:real_sru_split}
  \begin{tabular*}{\columnwidth}{@{\extracolsep{\fill}}lrrrr@{}}
    \toprule
    \textbf{Encoder} & \textbf{$n_{\mathrm{test}}$} &
    \textbf{Known} & \textbf{Unknown UNR} & \textbf{HC@.9} \\
    \midrule
    Qwen3-Emb    & 77 & .65 & .21 & .00 \\
    Jina v5      & 77 & .65 & .15 & .00 \\
    BioLORD      & 77 & .74 & .18 & .13 \\
    TF--IDF      & 77 & .70 & .18 & .13 \\
    \bottomrule
  \end{tabular*}
\end{table}

\section{Label-Withholding on Natural Clinical Corpora}
\label{app:label_withholding}

This bridge is \emph{not} counterfactual evidence deletion from natural notes.
We use Qwen3-14B-labeled MIMIC struct rows: training includes only clean
observed-positive/negative pairs (present vs.\ absent documentation in the
full note), while the held-out test split retains naturally
\emph{unmentioned} concept--note cells ($n_{\mathrm{unmentioned}}=1202$).
No evidence spans are removed at inference time; Unknown cases are those
Qwen3-14B marked \emph{unmentioned} in the discharge summary.
Because observation-label precision on $O^-$ is $\approx$.60
(Table~\ref{tab:audit}), treat this as a noisy-label sensitivity bridge
between controlled omission and real-note human gold rather than a gold
naturally occurring Unknown benchmark. Table~\ref{tab:label_withholding}
reports results on the matched MIMIC cohort. UNR remains high ($\geq$.75 for
Qwen3-Emb, Jina v5, BioLORD), confirming collapse on natural note silence
under binary training even without synthetic backgrounds. MedNLI evaluation
uses the same protocol when PhysioNet files are placed under
\texttt{data/mednli/}; we omit it here pending credential access.

\begin{table}[!htbp]
  \centering
  \scriptsize
  \renewcommand{\arraystretch}{1.06}
  \caption{MIMIC label-withholding bridge (train observed; test unmentioned).}
  \label{tab:label_withholding}
  \begin{tabular*}{\columnwidth}{@{\extracolsep{\fill}}lrrrr@{}}
    \toprule
    \textbf{Encoder} & \textbf{Known} & \textbf{UNR} &
    \textbf{HC@.9} & \textbf{Det.} \\
    \midrule
    Qwen3-Emb    & .81 & .80 & .72 & .48 \\
    Jina v5      & .79 & .81 & .73 & .46 \\
    BioLORD      & .72 & .75 & .67 & .45 \\
    TF--IDF      & .86 & .84 & .76 & .49 \\
    \bottomrule
  \end{tabular*}
\end{table}

\section{Simple Artifact Baselines}
\label{app:artifact_baselines}

Table~\ref{tab:artifact_baselines} compares simple artifact detectors against
dense encoders on the SilentEHR new-background protocol. Length, concept-string
match, and bag-of-words coverage either collapse to chance Known accuracy
(.50) or fail to reproduce high-confidence CWE. A lexical deny/affirm cue
classifier achieves high UNR (.87) but with low Unmentioned-detection AUROC
(.29), indicating a different failure mode than representation-driven collapse.
A random-embedding control matches chance UNR ($\approx$.50) with HCUNR@0.9
$=.37$, whereas Qwen3-Emb, Jina v5, and BioLORD retain HCUNR@0.9$\geq$.81
with Known$\geq$.96---supporting a representation-level rather than purely
lexical explanation of CWE on clean synthetic silence.

\begin{table}[!htbp]
  \centering
  \scriptsize
  \renewcommand{\arraystretch}{1.06}
  \caption{Artifact baselines vs.\ encoders (SilentEHR new-background).}
  \label{tab:artifact_baselines}
  \begin{tabular*}{\columnwidth}{@{\extracolsep{\fill}}lrrrr@{}}
    \toprule
    \textbf{Method} & \textbf{Known} & \textbf{UNR} &
    \textbf{HC@.9} & \textbf{Det.} \\
    \midrule
    Length           & .50 & .61 & .00 & .41 \\
    Concept match    & .50 & .00 & .00 & .50 \\
    Lexical cue      & .84 & .87 & .00 & .29 \\
    BoW coverage     & .50 & .00 & .00 & .50 \\
    Random encoder   & .48 & .50 & .37 & .47 \\
    \midrule
    Qwen3-Emb        & .98 & .81 & .86 & .67 \\
    Jina v5          & .96 & .60 & .81 & .50 \\
    BioLORD          & .99 & .80 & .99 & .18 \\
    \bottomrule
  \end{tabular*}
\end{table}

\section{Matched-Family Verbatim Example}
\label{app:matched_family}

Background~0 is shared across all nine realizations below
(\texttt{background\_id=0}, \texttt{concept\_id=penicillin\_allergy}).
Only the final sentence (or its absence) changes; the proposition remains
``The patient has a penicillin allergy.'' Unmentioned deletes the concept
sentence entirely rather than inserting a neutral filler (length matching is a
separate tier). Refuted states assert explicit absence; Not-assessed and
Uncertain retain a penicillin mention but do not assert present allergy.

\paragraph{Shared background (all states).}
\emph{The patient is a 67-year-old woman admitted with fatigue and poor oral
intake. Her medical history includes hypertension and hyperlipidemia. She was
hemodynamically stable on admission. Physical examination was otherwise
unremarkable. Labs showed mild hyponatremia. She lives with her daughter and is
independent in most ADLs. Home medications include amlodipine and atorvastatin.
Nursing notes describe improved appetite after IV fluids. Plan includes further
workup for failure to thrive.}

\paragraph{State-specific endings (verbatim from \texttt{controlled\_full.jsonl}).}
\begin{itemize}
  \item \textbf{Supported} (\texttt{p01440}): \emph{A penicillin allergy is documented.}
  \item \textbf{Supported, held template} (\texttt{p01441}): \emph{The record confirms penicillin allergy.}
  \item \textbf{Refuted} (\texttt{p01442}): \emph{The record states no known penicillin allergy.}
  \item \textbf{Refuted, held template} (\texttt{p01443}): \emph{Penicillin allergy is explicitly absent.}
  \item \textbf{Unmentioned} (\texttt{p01444}): \emph{(no additional sentence)}
  \item \textbf{Not assessed} (\texttt{p01445}): \emph{Penicillin allergy status was not documented.}
  \item \textbf{Not assessed, held} (\texttt{p01446}): \emph{Penicillin allergy reconciliation was deferred.}
  \item \textbf{Uncertain} (\texttt{p01447}): \emph{The reported penicillin allergy could not be confirmed.}
  \item \textbf{Uncertain, held} (\texttt{p01448}): \emph{The possible penicillin allergy requires clarification.}
\end{itemize}

\section{Human Annotation Protocol}
\label{app:annotation_protocol}

\paragraph{Sampling.}
We draw 300 admission--concept pairs from the Qwen3-14B-labeled MIMIC struct
bridge: 150 anemia + 150 thrombocytopenia. Stratification is by
(\emph{concept}, \emph{automatic mention status}, \emph{structured\_positive}),
with $n=25$ per cell. Mention status uses the struct-bridge coarse classes
\emph{observed\_pos}, \emph{observed\_neg}, and \emph{unmentioned} (six cells
per concept). This balances structured positive/negative laboratory state and
automatic mention class; we do \emph{not} perform additional Qwen-status
oversampling beyond this fixed grid, so Unknown is represented at the designed
cell rate rather than deliberately enriched.

\paragraph{Blind protocol.}
Annotators receive \texttt{mimic\_audit\_300.csv} with the full discharge note
(\texttt{encoder\_visible\_text}), concept name, proposition, structured
laboratory value (\texttt{struct\_value}), and chart/discharge timestamps.
They do \emph{not} see the automatic mention status, Qwen3-14B label,
four-cell \texttt{automatic\_cell}, or any model prediction. A separate key
file (\texttt{mimic\_audit\_300\_key.csv}) holds automatic labels for scoring
only. Annotators quote a short evidence span for every label other than
\emph{unmentioned}.

\paragraph{Seven-class definitions.}
\begin{itemize}
  \item \textbf{observed\_positive}: near-discharge text supports the abnormal
  laboratory state for the proposition.
  \item \textbf{observed\_negative}: near-discharge text supports normal/absent
  state.
  \item \textbf{unmentioned}: no usable textual evidence for the proposition.
  \item \textbf{admission-only}: evidence is explicitly limited to
  admission/early stay, not near discharge.
  \item \textbf{historical-only}: evidence is historical or resolved without a
  usable near-discharge state.
  \item \textbf{uncertain}: current evidence is explicitly uncertain or
  incomplete.
  \item \textbf{temporal-conflict}: incompatible near-discharge evidence cannot
  be reconciled.
\end{itemize}
Definitions match the Qwen3-14B observation prompt
(Appendix~\ref{app:observation_labeling}).

\paragraph{Consensus.}
Gold labels require exact agreement between both annotators on the seven-class
status. The original audit yields 255/300 consensus pairs; 45 disagreements are
\emph{excluded} from human gold (no third adjudicator). A later 50-pair
non-laboratory expansion on the same admissions reached 50/50 consensus and
was merged into \texttt{real\_clinical\_sru.jsonl} ($n=305$ consensus total).

\section{Real-Note Data Flow}
\label{app:data_flow}

Table~\ref{tab:data_flow} links the headline counts. The 493-pair silver set
is \emph{not} a superset of human gold: 200 struct pairs are sampled from
non-audit admissions (anemia/thrombocytopenia only), and 293 pairs expand audit
admissions with additional Qwen3-14B-labeled concepts on the same notes
(254/293 share an \texttt{audit\_id} with human rows). Consensus filtering
for S/R-only training removes Unknown labels from the \emph{training fold} only;
evaluation retains all Unknown rows, so label distribution at test time matches
the consensus set (159 Unknown / 146 known S+R in 5-fold CV).

\begin{table}[!htbp]
  \centering
  \scriptsize
  \renewcommand{\arraystretch}{1.08}
  \caption{Real-note evaluation data flow.}
  \label{tab:data_flow}
  \begin{tabular*}{\columnwidth}{@{\extracolsep{\fill}}llp{0.42\columnwidth}@{}}
    \toprule
    \textbf{Stage} & \textbf{$N$} & \textbf{Purpose} \\
    \midrule
    Qwen3-14B labeling pool & 13{,}799 notes & Observation statuses for MIMIC struct bridge \\
    MIMIC struct bridge & 38{,}043 pairs & All lab$\times$note cells with $C$ and $O$ \\
    Silver S/R/U subset & 493 pairs & Large-scale sensitivity (200 struct + 293 audit-note expansion) \\
    Human audit sample & 300 pairs & Dual blind annotation (150 anemia + 150 thrombocytopenia) \\
    Dual-annotator consensus & 305 pairs & 255 audit + 50 expansion; exact agreement only \\
    5-fold CV (primary) & 305 pairs & Pooled OOF: 146 S/R train labels, 159 Unknown evaluated \\
    70/30 sensitivity & 77 test pairs & Lab-audit consensus only (43 S/R + 34 Unknown); seed~0 \\
    \bottomrule
  \end{tabular*}
\end{table}

\section{Statistical Inference}
\label{app:statistics}

\paragraph{SilentEHR (family bootstrap).}
Matched variants share a background--concept family; we resample families with
replacement (2{,}000 replicates) and recompute UNR, known accuracy, and
HCUNR@0.9. Example new-background Unmentioned UNR 95\% intervals: Qwen3-Emb
[.90,.98], Jina v5 [.85,.93], BioLORD [.98,1.00]. Mitigation curves
(Table~\ref{tab:split_stability}; Appendix~\ref{app:split_stability}) report mean $\pm$ SD
over ten background redraws for UNR at 10/50/100 Unknown semantic groups.

\paragraph{Real-note gold (Wilson).}
Human-consensus S/R/U evaluation reports Wilson score 95\% intervals on pooled
out-of-fold binomial rates (Table~\ref{tab:real_sru_cv}). Audit-subset UPFNR
uses the same Wilson intervals (Table~\ref{tab:audit_sensitivity}).

\paragraph{Directional 0.5 reference.}
UNR$=0.5$ indicates equal Refuted vs.\ Supported pressure on Unknown profiles,
not chance performance on a three-way task. When multiple encoders exceed 0.5
with family-bootstrap lower bounds above 0.5, we interpret systematic Refuted
extrapolation rather than symmetric uncertainty.

\section{Observation-Label Inference Configuration}
\label{app:observation_labeling}

MIMIC observation labels were generated with \textbf{Qwen/Qwen3-14B-Instruct}
(instruct variant; thinking mode disabled) using local \textbf{vLLM 0.6.6}
inference. Decoding: \texttt{temperature=0}, \texttt{top\_p=1.0},
\texttt{max\_tokens=768}, \texttt{dtype=bfloat16}. One request per unique
admission (\texttt{hadm\_id}) batches all concepts for that note; discharge
text is passed without truncation (longest notes $\approx$48k characters).
Up to five retries with exponential backoff (cap 16\,s) on parse or transport
failures; persistent failures were dropped (zero failures in the archived run
of 13{,}799 notes).

\paragraph{Prompt (full text).}
\begin{quote}\small\ttfamily
You label observation status in de-identified clinical discharge-note text.
Use only the supplied TEXT; do not infer from outside knowledge or assume
that silence is negative.

The proposition concerns the patient's state near discharge. For each requested
concept, return exactly one label:
observed\_positive, observed\_negative, uncertain, historical\_only,
admission\_only, temporal\_conflict, or unmentioned.

Numeric thresholds: anemia = hemoglobin/Hgb $<$ 10; thrombocytopenia =
platelets/Plt $<$ 100; hyperkalemia = potassium/K $>$ 5.5. Return JSON only:
\{"labels":[\{"concept":"anemia","observation\_status":"unmentioned",
"temporal\_scope":"none","evidence\_span":"","implicit\_evidence":false\}]\}.
Evidence span must be an exact quote from TEXT, or empty when no evidence.
\end{quote}
The complete prompt with field definitions is in
\texttt{silentehr.py::OBSERVATION\_PROMPT} and
\texttt{configs/observation\_labeling.yaml}.

\paragraph{Parser.}
Responses must be a JSON object with a \texttt{labels[]} array. Each entry
requires \texttt{concept}, \texttt{observation\_status}, \texttt{temporal\_scope},
\texttt{evidence\_span}, and \texttt{implicit\_evidence}. Allowed statuses match
the seven-class schema in Appendix~\ref{app:annotation_protocol}.
\texttt{evidence\_span} must be an exact or whitespace-normalized substring of
the note; invalid spans are cleared and marked unverified. GPU model and count
were not retained in experiment logs; checkpoint name, prompt, decoding, and
parsing protocol above are sufficient to reproduce labeling behavior.
Restricted artifacts (raw label JSON with spans and note text) must not be
placed in a public repository.

\end{document}